\section{Fixed classical strength: locality and all-orders comparison}\label{sec:locality}
Let \[
 T=\begin{pmatrix}3&2\\4&3\end{pmatrix},\qquad M=T^{\mathsf T},\qquad
 \mathcal Cg=g\circ M.
\] On the two-torus choose a real even smooth function \(g\) with \[
 \operatorname{supp}g\cap\operatorname{supp}(\mathcal Cg)=\varnothing,
 \qquad f=g-\mathcal Cg.                                      \tag{LC.1}
\] There is an infinite-dimensional real vector space of such generators, constructed in Subsection~\ref{sec:locality}.2, containing nonzero functions. Let \(N\) be odd, let \(U_N\) be the cat quantization with the exact covariance below, and write \[
 G_N=\operatorname{Op}_N(g),\quad D_N=\operatorname{Op}_N(\mathcal Cg),
 \quad A_{{\kappa_0},N}=e^{i{\kappa_0} NG_N},\quad
 V_{{\kappa_0},N}=e^{i{\kappa_0} N(G_N-D_N)}U_N.                       \tag{LC.2}
\] The real parameter \({\kappa_0}\) is fixed as \(N\to\infty\). With the conventions below it represents a classical Hamiltonian time \({\kappa_0}/(2\pi)\), not a classical perturbation that shrinks with \(N\).

\textbf{Theorem A (all-orders comparison).} For each fixed bounded set of \({\kappa_0}\) values and each \(R>0\), \[
 \delta_N:=\|V_{{\kappa_0},N}-A_{{\kappa_0},N}U_NA_{{\kappa_0},N}^*\|_{\rm op}
       =O_{g,{\kappa_0},R}(N^{-R}).                              \tag{LC.3}
\] For every fixed \(p>0\) the same all-orders estimate holds uniformly after raising both unitaries to any integer power \(|t|\leq N^p\). For the explicit Gevrey-2 bump in Subsection~\ref{sec:locality}.2, constants \(c,C>0\) also give \[
 \delta_N\leq C{\kappa_0}^2N^2e^{-cN^{1/3}}                      \tag{LC.4}
\] for sufficiently large \(N\). This conservative exponent is not claimed optimal. It gives vanishing error even for \(|t|\leq e^{cN^{1/3}/2}\).

\textbf{Theorem B (classical interpretation).} Let \(\Phi_g^t\) be the Hamiltonian flow specified by \(\frac{d}{dt}(a\circ\Phi_g^t)=\{g,a\}\circ\Phi_g^t\), with \(\{g,a\}=g_q a_p-g_p a_q\). The map quantized by (LC.2) is \[
 M_{\kappa_0}=M\circ\Phi_f^{{\kappa_0}/(2\pi)}
 =\Phi_g^{-{\kappa_0}/(2\pi)}\circ M\circ\Phi_g^{{\kappa_0}/(2\pi)}. \tag{LC.5}
\] It is Anosov and mixing for area measure. For any nonzero example \(g\) and all sufficiently small fixed \({\kappa_0}\ne0\), it is nonlinear and differs from \(M\). Its conjugacy to \(M\) is essential, not a symmetry that may be omitted when interpreting its spectral statistics.

\textbf{Theorem C (microscopic stability along Pell moduli).} Let \(K\) tend to infinity through odd primes, put \(N+F\sqrt2=(1+\sqrt2)^K\), and restrict to the odd parity space of dimension \(D=(N-1)/2=Kn\). For fixed \({\kappa_0}\) and sufficiently large \(N\), the spectrum consists of \(K\) clusters, each containing exactly \(n\) eigenvalues, whose angular halfwidth \(a_N\) is \(O(N^{-R})\) for every \(R\). After a common rotation their centres are the \(K\)th roots of unity.

Under unfolding by the mean density \(D\), the spacing measure tends weakly to \(\delta_0\). More quantitatively, for every fixed finite \(p\geq1\), \[
 W_p\left(\mu_{\rm gaps},(1-1/n)\delta_0+(1/n)\delta_n\right)
       \leq2Da_N.                                           \tag{LC.6}
\] For each fixed \(x>0\), eventually the normalized ordered close-pair count is \textbf{exactly} \(n-1\). The local point process viewed from a uniformly random origin converges vaguely in probability to the empty process, despite having mean intensity one at every finite \(N\).

Spacing collapse, divergence of close-pair counts, and this empty local limit remain true under any increasing degree-one unfolding \(F_N\) satisfying \(\|F_N'\|_\infty\leq C N^A\) for fixed \(C,A\). Thus the present conclusion cannot be removed by an unfolding with polynomially bounded distortion. No assertion is made about unfoldings whose derivatives grow faster than every power of \(N\).

\subsection{Quantization and a finite locality estimate}\label{locality:quantization-and-a-finite-locality-estimate}

For an integer frequency \(v=(a,b)\) use the \textbf{physical integer Fourier lift} \[
 (Q_N(a,b)u)(x)=e^{2\pi ia(x-b/2)/N}u(x-b\bmod N).
\] The half in this exponent is ordinary division, not the inverse of two in \(\mathbb Z/N\mathbb Z\). For the modular-centred convention \(W_N\) used in the earlier notes, \(Q_N(a,b)=(-1)^{ab}W_N(a,b)\). In particular, \[
 Q_N(v)Q_N(w)=e^{\pi i\omega(v,w)/N}Q_N(v+w),\qquad
 \omega(v,w)=v_1w_2-v_2w_1.                                 \tag{LC.7}
\] Each \(Q_N(v)\) is unitary, \(Q_N(v)^*=Q_N(-v)\), and parity sends \(Q_N(v)\) to \(Q_N(-v)\). Set \[
 \operatorname{Op}_N(a)=\sum_{v\in\mathbb Z^2}\widehat a(v)Q_N(v),
 \quad \mathcal M_r(a)=\sum_v|\widehat a(v)|(1+|v|_1)^r.
\] The series converges absolutely in operator norm for smooth \(a\); \(\|\operatorname{Op}_N(a)\|\leq\mathcal M_0(a)\), and real symbols give Hermitian matrices. Frequency aliases cause no problem in this definition.

Since \(T\) is the identity modulo two, its finite Weyl implementer also has exact physical covariance \[
 U_NQ_N(v)U_N^*=Q_N(Tv),\qquad
 U_N\operatorname{Op}_N(a)U_N^*=\operatorname{Op}_N(a\circ M). \tag{LC.8}
\] One such implementer has the kernel \[
 U_N(x,y)=\frac{c_N}{\sqrt N}
 e_N\!\left((3x^2-2xy+3y^2)/8\right),\quad |c_N|=1,         \tag{LC.9}
\] where division by eight is modular. Orthogonality of additive characters proves unitarity; direct substitution on the two Weyl generators proves (LC.8). The kernel is parity-preserving. Its global phase can be fixed later.

\textbf{Lemma 1 (disjoint-symbol product).} If \(a,b\in C^\infty(\mathbb T^2)\) have disjoint closed supports, then for every integer \(r\geq1\), \[
 \|\operatorname{Op}_N(a)\operatorname{Op}_N(b)\|_{\rm op}
 \leq\frac{\pi^r}{r!N^r}\mathcal M_r(a)\mathcal M_r(b).       \tag{LC.10}
\]

\textbf{Proof.} Expand (LC.7) and Taylor-expand \(e^{i\pi\omega/N}\) through degree \(r-1\). For real \(z\) the integral Taylor remainder satisfies \(|e^{iz}-\sum_{j<r}(iz)^j/j!|\leq|z|^r/r!\), with no restriction on \(z\). For each \(j<r\) the coefficient sum \[
 \sum_{v,w}\widehat a(v)\widehat b(w)\omega(v,w)^j Q_N(v+w)
\] is the quantization of \[
 \frac{1}{(-4\pi^2)^j}\sum_{l=0}^j(-1)^l\binom jl
 (\partial_q^{j-l}\partial_p^l a)
 (\partial_q^l\partial_p^{j-l} b).                           \tag{LC.11}
\] All these functions vanish, including \(ab\) at \(j=0\), because derivatives are supported in the respective original closed supports. Their Fourier cancellation occurs before applying \(Q_N\), so aliases do not invalidate it. The remaining absolutely convergent sum is bounded using \(|\omega(v,w)|\leq|v|_1|w|_1\) and \(\|Q_N(v+w)\|=1\). This proves (LC.10).

For real \(a,b\), their two products are adjoints, so \[
 \|[\operatorname{Op}_N(a),\operatorname{Op}_N(b)]\|
 \leq\frac{2\pi^r}{r!N^r}\mathcal M_r(a)\mathcal M_r(b).      \tag{LC.12}
\] This all-orders statement is stronger than just Poisson commutation.

\subsection{An explicit infinite-dimensional family and a Gevrey bound}\label{locality:an-explicit-infinite-dimensional-family-and-a-gevrey-bound}

Put \(r_0=1/16\) and let \(S\) be the union of the closed torus squares of halfwidth \(r_0\) centred at \((1/4,0)\) and \((3/4,0)\). If \((q,p)\in S\), the second coordinate \(2q+3p\) of \(M(q,p)\) has distance at least \[
 1/2-5r_0=3/16>r_0                                          \tag{LC.13}
\] from the integers. Therefore \(S\cap M^{-1}S=\varnothing\). Every real even smooth function supported in \(S\) satisfies (LC.1). This is an infinite-dimensional space. For nonzero \(g\), \(f\) is nonzero because it equals \(g\) on \(S\). Also \(\int f=0\), regardless of the mean of \(g\).

Here is the fixed example used numerically. Write \(\langle x\rangle\) for the representative of \(x\bmod1\) in \([-1/2,1/2)\), set \[
 \beta(u)=\begin{cases}\exp(1-1/(1-u^2)),&|u|<1,\\0,&|u|\geq1,\end{cases}
 \quad u_\pm(q)=\langle q\mp1/4\rangle/r_0,
\] and choose \(\eta=1/3\). Then \[
 \begin{split}
 A(q)&=\beta(u_+(q))(1+\eta u_+(q))
                +\beta(u_-(q))(1-\eta u_-(q)),\\
 B(p)&=\beta(\langle p\rangle/r_0),\qquad g(q,p)=A(q)B(p).
 \end{split}                                                \tag{LC.14}
\] This \(g\) is real, even and nonzero. The unequal tilts of its two lobes avoid a spurious half-period translation symmetry in the numerical example; no absence of all other symmetries is claimed.

For completeness, its Gevrey regularity and the nonoptimal stretched exponential bound need no black-box pseudodifferential theorem. If \(D_t=1-t^2>0\), use the complex disk \(|z-t|\leq D_t/8\). On this disk \[
 |(1-z^2)-D_t|\leq17D_t/64,\qquad
 \mathop{\rm Re}(1-z^2)^{-1}\geq\frac{3008}{6561D_t}>
 \frac1{3D_t}.
\] Cauchy's estimate and maximization of \(D_t^{-j}e^{-1/(3D_t)}\) give \[
 |\beta^{(j)}(t)|\leq e\,24^j(j!)^2.                         \tag{LC.15}
\] All derivatives extend by zero at \(\pm1\). Multiplication by \(1\pm\eta t\) changes the prefactor to at most \(2e\), with the same base 24.

A concrete Fourier-moment estimate follows. Put \(L_0=24/r_0\), \(B_0=9L_0/\pi\) and \[
 A_0=8r_0+(8/3)r_0eL_0^2,\qquad
 B_1=2r_0+(2/3)r_0eL_0^2.
\] Integrating the one-dimensional Fourier coefficients by parts \(j+2\) times and using \(\sum_{k\ne0}k^{-2}=\pi^2/3\) proves \[
 \sum_k(1+|k|)^j|\widehat A(k)|\leq A_0B_0^j(j!)^2,
 \quad
 \sum_k(1+|k|)^j|\widehat B(k)|\leq B_1B_0^j(j!)^2.           \tag{LC.16}
\] One uses \((j+2)^2(j+1)^2\leq4\,9^j\). Expanding the sum of the two frequency weights gives \[
 \mathcal M_j(g)\leq A_0B_1(2B_0)^j(j!)^2,
 \quad
 \mathcal M_j(\mathcal Cg)\leq A_0B_1(14B_0)^j(j!)^2,        \tag{LC.17}
\] since \(\|T\|_1=7\). These explicit constants are intentionally coarse and are not useful finite-size error bars for the numerical matrices below.

More generally, suppose \(\mathcal M_j(g)\leq A_gB_g^j(j!)^s\) and \(\mathcal M_j(\mathcal Cg)\leq A_dB_d^j(j!)^s\) for some \(s>1\). The comparison proved in Subsection~\ref{sec:locality}.3 has the bound \[
 \delta_N\leq{\kappa_0}^2N^2A_gA_d(C/N)^j(j!)^{2s-1},
 \qquad C=\pi B_gB_d.                                      \tag{LC.18}
\] For \(j=\lfloor(N/(2C))^{1/(2s-1)}\rfloor\geq1\), use \(j!\leq j^j\) to bound the last factor by \(2^{-j}\). This proves a bound \(C'{\kappa_0}^2N^2\exp(-cN^{1/(2s-1)})\), and (LC.4) follows for \(s=2\). The general factorial-growth/optimal-truncation mechanism is established Gevrey analysis; this section does not claim to originate it or to improve the best Gevrey remainders in other quantization settings.

\subsection{Fixed-strength conjugacy, classical dynamics, and long times}\label{locality:fixed-strength-conjugacy-classical-dynamics-and-long-times}

For Hermitian \(H,J\) and real \(t\), the two-factor Duhamel inequality gives \[
 \|e^{it(H+J)}-e^{itH}e^{itJ}\|\leq(t^2/2)\|[H,J]\|.        \tag{LC.19}
\] One proof differentiates \(e^{i(t-u)(H+J)}e^{iuH}e^{iuJ}\) and integrates the conjugation difference of \(J\) twice. Unitary invariance of the norm leaves a triangular integration domain of area \(t^2/2\). In particular, \(|t|\|H\|\) need not be small.

By (LC.8), \[
 A_{{\kappa_0},N}U_NA_{{\kappa_0},N}^*
       =e^{i{\kappa_0} NG_N}e^{-i{\kappa_0} ND_N}U_N.
\] Apply (LC.19) with \(t={\kappa_0} N\), \(H=G_N\), \(J=-D_N\), and then (LC.12). For every integer \(r\geq1\), \[
 \boxed{\displaystyle
 \delta_N\leq\frac{{\kappa_0}^2\pi^r}{r!}N^{2-r}
                   \mathcal M_r(g)\mathcal M_r(\mathcal Cg).} \tag{LC.20}
\] This proves Theorem A and (LC.18). Notice that we have not set the finite commutator to zero, nor silently replaced the one-exponential kick by a product of two kicks.

For integer \(t\), telescoping products of unitaries, or their adjoints when \(t<0\), yields \[
 \|V_{{\kappa_0},N}^t-A_{{\kappa_0},N}U_N^tA_{{\kappa_0},N}^*\|
       \leq |t|\delta_N.                                   \tag{LC.21}
\] Choosing \(r\) larger than any fixed polynomial time exponent proves the all-polynomial-time assertion. The same estimate proves the stretched exponential-time assertion from (LC.4). Normalized traces differ by at most \(|t|\delta_N\); their absolute squares differ by at most \(2|t|\delta_N\). These statements control Heisenberg times \(t\asymp D\), not just \(O(\log N)\).

To establish (LC.5), note that (LC.1) implies \(\{g,\mathcal Cg\}=0\), so the two Hamiltonian flows commute and \[
 \Phi_{\mathcal Cg}^t=M^{-1}\Phi_g^tM,\qquad
 \Phi_f^t=\Phi_g^t\Phi_{\mathcal Cg}^{-t}.
\] Consequently \(M\Phi_f^t=\Phi_g^{-t}M\Phi_g^t\). This is an exact classical identity for every finite \(t\). A smooth symplectic conjugacy transports the invariant hyperbolic splitting, with bounded derivatives in both directions, so it preserves uniform hyperbolicity. It also transports area measure; mixing follows from mixing of \(M\). The latter can be checked on Fourier characters: for nonzero integer \(v\), \(|T^nv|\to\infty\), and hence fixed Fourier modes eventually have zero correlation. Density extends the result to square-integrable observables.

The normalization of classical time is checked rather than presumed: (LC.7) gives \[
 [\operatorname{Op}_N(g),\operatorname{Op}_N(a)]
 =-\frac{i}{2\pi N}\operatorname{Op}_N(\{g,a\})+O(N^{-3})
\] for fixed smooth symbols. Differentiating unitary conjugation and using Duhamel on a fixed compact time interval gives the pullback by \(\Phi_g^{{\kappa_0}/(2\pi)}\), with a remainder tending to zero. Thus (LC.2) is a quantization of the fixed map in (LC.5).

For nonzero \(g\) the smooth function \(f\) is nonconstant. Its flow moves some point for every sufficiently small nonzero time. The complement of the support of \(f\) has nonempty interior (the support occupies less than the whole torus), and the flow is exactly the identity there. Accordingly \(M_{\kappa_0}\) equals \(M\) on an open set but is not equal to it everywhere for small nonzero \({\kappa_0}\). An affine torus map agreeing with \(M\) on an open set agrees everywhere. Therefore these fixed maps are genuinely nonlinear, although smoothly conjugate to the same linear map.

\subsection{Arithmetic multiplicities and individual cluster control}\label{locality:arithmetic-multiplicities-and-individual-cluster-control}

The following finite arithmetic facts can be established without primality of \(N\). Write \(\alpha=1+\sqrt2\), \(\lambda=\alpha^2\) and \(N+F\sqrt2=\alpha^K\) with \(K\) an odd prime. The Pell identity gives \(N^2-2F^2=-1\), hence \(F\) is invertible modulo \(N\). Set \[
 s=F^{-1}\pmod N,\quad s^2=2,\quad h=3+2s.
\] Since \((1+s)^K=N+Fs=1\), \(h^K=1\). Also \(h-1=2(1+s)\) is a unit. Modulo each prime divisor of \(N\), the order of \(h\) is exactly \(K\). Thus \(1-h^j\) is a unit for \(1\leq j<K\), and \(1+h^j\) is a unit for every \(j\).

The symplectic matrix \[
 P=\begin{pmatrix}1&-1/(2s)\\s&1/2\end{pmatrix}
\] satisfies \(T=P\operatorname{diag}(h,h^{-1})P^{-1}\). Its lower-left entry is a unit, so its finite modular Weyl implementer has the same type of quadratic kernel as (LC.9) and preserves parity. Conjugating \(S_h^{-1}\) by this implementer constructs a phase-normalized \(U_N\) with \(U_N^K=I\). Here \(S_hu(x)=u(h^{-1}x)\). The unit conditions above show that every nonzero residue has an orbit of length \(K\), and that an orbit and its negative are distinct. On the odd space, there are \(n=(N-1)/(2K)\) cyclic orbits, each contributing every \(K\)th root once. Hence the multiplicity is exactly \(n\) for each root. A different global phase merely rotates all centres and changes none of the spacing conclusions.

Let \(W=A_{{\kappa_0},N}U_NA_{{\kappa_0},N}^*\) on the odd space. Restriction cannot increase (LC.20). For large \(N\) take \(\delta_N<\sin(\pi/K)\). Spectral stability for a normal matrix places each eigenvalue of \(V\) within \(\delta_N\) of a \(K\)th root. To verify multiplicities, let \(E=VW^*\) and use its principal Hermitian logarithm, \(E=e^{iH}\). Then \(\|H\|\leq2\arcsin(\delta_N/2)\), and the path \(e^{iuH}W\), \(0\leq u\leq1\), stays within \(\delta_N\) of \(W\). The disjoint slightly enlarged root disks have constant Riesz ranks along this path. Thus every disk contains exactly \(n\) eigenvalues. Their circle-coordinate halfwidth is bounded by \[
 a_N=\pi^{-1}\arcsin(\delta_N/2)=O(N^{-\infty}).              \tag{LC.22}
\] This is an operator-norm, individual-cluster statement, unlike the earlier normalized Hilbert--Schmidt comparison. It still makes no claim about individual eigenvectors inside a nearly degenerate cluster.

\subsection{Gaps, local counts, and robustness against unfolding}\label{locality:gaps-local-counts-and-robustness-against-unfolding}

Sort the \(D\) eigenangles cyclically, retaining multiplicities, and put \(S_l=D(\theta_{l+1}-\theta_l)\), with the cyclic final gap. Exactly \(D-K\) gaps are within clusters and at most \(2Da_N\). The other \(K\) gaps are in \([n-2Da_N,n+2Da_N]\). Matching these two classes to zeros and copies of \(n\) proves (LC.6). Since \(1/n=K/D\to0\), the reference measures in (LC.6) converge weakly to \(\delta_0\), and so do the actual gap measures. Their means are nevertheless exactly one at each \(N\): rare large gaps carry that mean.

Define the ordered close-pair count \[
 R_N(x)=\frac1D\#\{(l,j):l\ne j,\ D\operatorname{dist}_{\mathbb T}
                         (\theta_l,\theta_j)\leq x\}.
\] For fixed \(x>0\), eventually \(2Da_N<x\) and \(D(1/K-2a_N)>x\). Every within-cluster ordered pair is counted and no cross-cluster pair is counted. Therefore \[
 \boxed{R_N(x)=n-1\quad\hbox{eventually}.}                   \tag{LC.23}
\] For a fixed compactly supported continuously differentiable test \(\phi\), its periodized factorial statistic likewise obeys \[
 \mathcal R_N(\phi)=(n-1)\phi(0)
       +O\bigl((n-1)Da_N\|\phi'\|_\infty\bigr).             \tag{LC.24}
\] The remainder is superpolynomially small. This formula is not asserted with this rate for an arbitrary continuous test.

Choose an origin uniformly on the circle. A window of unfolded length \(L\) can meet a cluster only if its origin belongs to one of \(K\) arcs of length at most \(L/D+2a_N\). Hence \[
 \mathbb P\{\hbox{window contains a level}\}
       \leq KL/D+2Ka_N\longrightarrow0.                     \tag{LC.25}
\] The same argument applies to any fixed compact observation window, giving vague convergence in probability of the translated local point measures to the empty measure. Their expected count is still \(L\); uniform integrability fails. This is a point-process statement, not an inference from diverging second moments alone.

Now let \(F_N:\mathbb R\to\mathbb R\) be increasing with \(F_N(x+1)=F_N(x)+1\) and Lipschitz constant at most \(C N^A\). Image clusters have angular diameter at most \(2CN^Aa_N\). There are at most \(K\) gaps not internal to these clusters. Internal unfolded gaps tend uniformly to zero, proving spacing collapse. All \(n(n-1)\) ordered pairs in each original cluster are still close, so \(R_N(x)\geq n-1\) eventually. Finally (LC.25) is replaced by \[
 \mathbb P\{\hbox{image window contains a level}\}
       \leq KL/D+2KCN^Aa_N\longrightarrow0.                 \tag{LC.26}
\] No lower bound on \(F_N'\) is used, and image clusters need not remain uniformly separated. Thus (LC.23) need not remain an equality after unfolding; the lower bound and empty-process conclusion are the correct statements.

\subsection{Flat directions, hidden structure, and a falsification control}\label{locality:flat-directions-hidden-structure-and-a-falsification-control}

For every generator in (LC.1), \(\{g,\mathcal Cg\}=0\). In the previous smooth notation, \(\Gamma(g)=\sigma^2(\{g,\mathcal Cg\}/(4\pi))=0\). Thus the second-order exceptional set contains an explicit infinite-dimensional linear subspace containing nonzero smooth generators. This does not contradict openness and density of \(\Gamma>0\); an infinite-dimensional exceptional subspace can still have empty interior. Subtracting the mean of \(g\) places it in the zero-mean class without changing \(f\), \(\Gamma\), or the quantum commutator. One should apply the support argument to the original bump, not falsely describe its mean-subtracted version as compactly supported.

Poisson commutation \textbf{alone} is insufficient for (LC.3). For a negative control take \(a=b_{(1,0)}+b_{(0,1)}\), where \(b_v=2\cos(2\pi v\cdot x)\), and \(b=a^2\). Of course \(\{a,b\}=0\). Yet the cubic Weyl coefficient \[
 \sum_{v+w=k}\widehat a(v)\widehat b(w)\omega(v,w)^3
\] is nonzero: at frequencies \((-2,-1),(-2,1),(-1,-2),(-1,2)\) its values are \(-6,6,6,-6\), respectively, and the four opposite frequencies give the same values. The commutator therefore has a nonzero term of order \(N^{-3}\). For a sufficiently large odd \(N\) these finite frequencies are distinct modulo \(N\), so this is not a formal nonzero coefficient that disappears under quantization. Both implementations retain this negative control. The all-orders cancellation in (LC.11), not just its first term, is essential.

Finally, if \(J_N\) commutes with \(U_N\) and \(\|J_N\|\leq1\), then \[
 \|[V_{{\kappa_0},N},A_{{\kappa_0},N}J_NA_{{\kappa_0},N}^*]\|
        \leq2\delta_N.                                    \tag{LC.27}
\] The arithmetic commutant is therefore transported to an all-orders approximate commutant. These systems do not evade arithmetic structure by being nonlinear. We do not claim their complete symmetry reduction, simple spectrum, or generic random-matrix behaviour. Obtaining corresponding positive universality for fixed maps not conjugate to arithmetic models is not accomplished by this construction.
